% Generated draft: requires mathematical and source review.
\documentclass[11pt]{article}
\usepackage[margin=25mm]{geometry}
\usepackage{fontspec}
\setmainfont{Latin Modern Roman}
\usepackage{amsmath,amssymb,mathtools,unicode-math}
\setmathfont{Latin Modern Math}
\usepackage{enumitem,needspace,xurl,hyperref}
\setlength{\parindent}{0pt}
\setlength{\parskip}{5pt}
\setlength{\emergencystretch}{3em}
\newcommand{\N}{\mathbb{N}}
\newcommand{\Z}{\mathbb{Z}}
\newcommand{\Q}{\mathbb{Q}}
\newcommand{\R}{\mathbb{R}}
\newcommand{\C}{\mathbb{C}}
\newcommand{\F}{\mathbb{F}}
\newcommand{\A}{\mathbb{A}}
\newcommand{\PP}{\mathbb{P}}
\newcommand{\E}{\mathbb{E}}
\newcommand{\eps}{\varepsilon}
\newcommand{\dd}{\,\mathrm{d}}
\begin{document}
\section*{UnsolvedMath ID: 30003095}
\textbf{Source problem code:} catalogue entry
\subsection{Definitions and mathematical statement}
The supplied problem asks for an efficient numerical method for finding a diffeomorphism that transports one smooth density to another. The underlying space and the admissible densities are not specified, so the following gives the transport equation conditionally, without asserting a universal existence or algorithmic theorem.

In a formulation on a specified smooth manifold or smooth domain \(M\), let \(\mu_0\) and \(\mu_1\) be smooth densities interpreted as measures. A diffeomorphism \(\varphi\colon M\to M\) transports \(\mu_0\) to \(\mu_1\) if
\[
\varphi_*\mu_0=\mu_1,
\qquad
(\varphi_*\mu_0)(B):=\mu_0\bigl(\varphi^{-1}(B)\bigr)
\]
for every measurable subset \(B\subseteq M\). Equivalently, \(\varphi^*\mu_1=\mu_0\). In particular, transport requires matching total masses, \(\mu_0(M)=\mu_1(M)\); this compatibility condition is not explicitly stated in the source.

\textbf{Numerical problem.} Given two smooth densities in an appropriate specified setting admitting such a diffeomorphism, develop an efficient numerical procedure for computing or approximating a diffeomorphism satisfying the transport equation.

This remains a research question rather than a fully quantified algorithmic assertion: the source does not specify the dimension or geometry of \(M\), compactness, boundary conditions, strict positivity versus possible vanishing of the densities, numerical input representation, approximation norm, error tolerance, convergence requirements, or a measure of computational efficiency. These choices, and hypotheses ensuring existence of the required diffeomorphism, remain unresolved qualifications of the formulation.
\subsection{Short English statement}
Develop an efficient numerical method to construct a smooth invertible map that carries one prescribed smooth density to another. The source leaves the geometric setting, existence hypotheses, boundary conditions, accuracy requirements, and meaning of efficiency unspecified.
\subsection{Sources}
\begin{itemize}
\item Local registry: \texttt{data/unsolvedmath/problems.json}, permanent UnsolvedMath ID \(30003095\), source problem code \texttt{catalogue entry}, title \emph{Efficient Numerical Transport between Smooth Densities}. The supplied record assigns the problem to 2016 and names no individual proposer.
\item Supplied source citation: \emph{Geometric Numerical Integration} (2016), Oberwolfach Reports, pp.\ 869--948, DOI \texttt{10.4171/owr/2016/18}. Supplied literature-source label: OWR 2016 report. \url{https://doi.org/10.4171/owr/2016/18}. The supplied background describes Oberwolfach Reports as documenting mathematical workshops and their open-problem discussions.
\item Supplied literature reference: Bauer--Joshi--Modin, DOI \texttt{10.1137/151006238}. \url{https://doi.org/10.1137/151006238}. No further bibliographic details are supplied.
\item Literature assessment dated exactly \texttt{2026-08-21}: the supplied curation labels the problem \emph{partially solved} and states that optimal information transport supplies efficient numerical algorithms for smooth density matching, while the broad numerical density-transport challenge is not solved in full generality. This is supplied dated curation derived from the OWR record, not an independently verified assessment. Its time-sensitive status and cited claims require independent verification.
\end{itemize}
\end{document}
